\section{Fundamentos teóricos}

\subsection{Señal BOLD y función de respuesta hemodinámica}

La señal BOLD depende de cambios en flujo sanguíneo, volumen vascular y concentración de desoxihemoglobina. La HRF resume la respuesta temporal observada después de una entrada neuronal breve. En un modelo lineal, la señal esperada puede expresarse como

\begin{equation}
    y(t) = (u * h)(t) + \varepsilon(t),
\end{equation}

donde $u(t)$ representa el estímulo, $h(t)$ la HRF y $\varepsilon(t)$ el término de error. El GLM utiliza esta relación para construir regresores y estimar amplitudes de activación \citep{Friston1998}.

\subsection{Modelo Balloon--Windkessel}

Se utilizaron cuatro estados normalizados: señal vasodilatadora $s(t)$, flujo de entrada $f(t)$, volumen sanguíneo venoso $v(t)$ y contenido de desoxihemoglobina $q(t)$. La dinámica implementada fue

\begin{align}
\dot{s}(t) &= \epsilon u(t) - \kappa_s s(t) - \kappa_f\left[f(t)-1\right], \\
\dot{f}(t) &= s(t), \\
\tau\dot{v}(t) &= f(t) - v(t)^{1/\alpha}, \\
\tau\dot{q}(t) &= \frac{f(t)}{E_0}\left[1-(1-E_0)^{1/f(t)}\right]
                  - v(t)^{1/\alpha}\frac{q(t)}{v(t)}.
\end{align}

La salida BOLD fraccional se calculó mediante

\begin{equation}
 y(t)=V_0\left[k_1(1-q(t))+k_2\left(1-\frac{q(t)}{v(t)}\right)+k_3(1-v(t))\right].
\end{equation}

El modelo captura transitorios como sobreimpulso y undershoot postestímulo, además de la interacción entre flujo, volumen y oxigenación \citep{Buxton1998,Friston2000}.

\subsection{Redes neuronales informadas por la física}

Una PINN aproxima los estados mediante una red neuronal diferenciable y penaliza el incumplimiento de las ecuaciones gobernantes. Dado un conjunto de tiempos de colocación $\{t_j\}$, la pérdida física se construye con residuos obtenidos por diferenciación automática. Esta estrategia permite abordar problemas directos e inversos sin discretizar explícitamente todas las variables latentes \citep{Raissi2019}.

En este trabajo, $\epsilon$, $\tau$ y $\alpha$ se trataron como parámetros entrenables acotados. Los restantes parámetros fisiológicos se mantuvieron fijos para reducir la no identificabilidad.
